\documentclass{article}
\usepackage[margin=1in]{geometry}
\usepackage{longtable}
\begin{document}

\section*{Phase 8 Task 2 - Frozen End-to-End Reference Path}

\textbf{Source of truth:} \texttt{Documentation/SRS.tex}, \texttt{Documentation/phases/phase8.tex}, and the generated freeze manifest under \texttt{reports/phase8/reference\_window\_freeze/}.\\
\textbf{Task goal:} Freeze one defendable end-to-end reproducibility path from raw archive through replay, candidate generation, alerting, validation, and ML or Phase 7 research packaging, with exact paths, versions, and hashes.\\
\textbf{Status in this repo snapshot:} The reproducibility path is now frozen \textbf{at the provenance level}. The exact code, docs, schema files, expected runtime artifact paths, and database table targets are all enumerated and hashable. However, the current local workspace does \textbf{not} contain materialized runtime rows or saved report artifacts for the later stages, so the frozen chain is presently a \textbf{defendable definition plus evidence-gap record}, not a fully populated runtime packet.

\section*{Chosen Reference Window}
\textbf{Window:} \texttt{2026-04-20T05:00:00+00:00} to \texttt{2026-04-20T06:00:00+00:00}\\
\textbf{Primary source system:} \texttt{gamma\_events}

\textbf{Why this window was chosen}
\begin{itemize}
\item It is the only exact one-hour raw/replay proof window explicitly named in a canonical committed phase document.
\item \texttt{Documentation/phases/phase2.tex} already records raw/archive parity and deterministic republish evidence for this hour.
\item Freezing a later anonymous window without committed proof would be weaker than freezing the exact committed one and recording what downstream stages still need to materialize.
\end{itemize}

\section*{Frozen Chain Definition}
\begin{longtable}{p{0.20\linewidth}p{0.34\linewidth}p{0.34\linewidth}}
\textbf{Stage} & \textbf{Frozen Artifacts} & \textbf{Expected Runtime Output} \\
Raw archive and replay & \texttt{Documentation/phases/phase2.tex}, \texttt{Documentation/phases/phase2\_gate2\_signoff.tex}, \texttt{utils/event\_log.py}, \texttt{validation/phase2\_replay.py}, \texttt{validation/phase2\_republish.py} & \texttt{data/raw/year=2026/month=04/day=20/hour=05/source\_system=gamma\_events/events.ndjson}, \texttt{data/detector\_input/...}, and the republished replay output path cited in the Phase 2 delivery doc \\
Candidate generation & \texttt{Documentation/phases/phase3.tex}, \texttt{database/PHASE3\_LOCAL\_RUNBOOK.md}, \texttt{phase3/detector.py}, \texttt{run\_phase3\_live.py} & persisted candidate and feature rows in \texttt{signal\_candidates}, \texttt{signal\_episodes}, and \texttt{signal\_features} \\
Alert and evidence & \texttt{Documentation/phases/phase4.tex}, \texttt{Documentation/phases/phase4\_gate4\_signoff.tex}, \texttt{phase4/evidence.py}, \texttt{phase4/alerts.py}, \texttt{run\_phase4\_pipeline.py} & persisted evidence and alert rows in \texttt{evidence\_snapshots}, \texttt{alerts}, and \texttt{alert\_delivery\_attempts} \\
Validation and backtest & \texttt{Documentation/phases/phase5.tex}, \texttt{phase5/replay.py}, \texttt{phase5/simulator.py}, \texttt{phase5/reporting.py}, \texttt{run\_phase5\_replay.py} & Phase 5 replay/validation/backtest outputs under \texttt{reports/phase5/} and stored summary rows in \texttt{validation\_runs} and \texttt{backtest\_artifacts} \\
ML / research endpoint & \texttt{phase6/training.py}, \texttt{run\_phase6\_train\_ranker.py}, \texttt{phase7/graph\_features.py}, \texttt{phase7/packaging.py}, \texttt{run\_phase7\_build\_research\_package.py} & a Phase 6 evaluation artifact or a Phase 7 research package rooted under \texttt{reports/phase7/} and backed by ledger rows in \texttt{phase7\_experiment\_ledger} \\
\end{longtable}

\section*{Generated Freeze Manifest}
Task 2 introduces one machine-readable manifest and one human-readable summary:
\begin{itemize}
\item \texttt{reports/phase8/reference\_window\_freeze/phase8\_reference\_window\_manifest.json}
\item \texttt{reports/phase8/reference\_window\_freeze/phase8\_reference\_window\_summary.md}
\end{itemize}

These files record:
\begin{itemize}
\item the chosen window and why it was selected,
\item the current git commit,
\item exact committed file paths and SHA-256 hashes for every frozen code and doc input,
\item exact runtime artifact paths expected for the same chain,
\item version constants for Phase 3 through Phase 7,
\item database table targets and current row counts,
\item and a stage-by-stage status saying whether the path is fully materialized or only frozen at the provenance level.
\end{itemize}

\section*{Current Evidence Gap}
The local committed SQLite file contains the later-phase schema, but the tables used by Phases 2 through 7 are empty in this workspace snapshot. The \texttt{reports/phase7/} directory exists, but it contains no committed research payloads. Therefore:
\begin{itemize}
\item the code path and document provenance are now frozen and hashable,
\item the exact runtime artifact targets are frozen and named,
\item but the actual end-to-end runtime evidence packet still needs to be materialized in a later closeout pass.
\end{itemize}

\section*{Why This Still Counts as Task 2 Progress}
Task 2 is about defending one end-to-end path with exact paths, versions, and hashes. This repo snapshot could not honestly claim a fully populated end-to-end packet, because the runtime outputs are absent. The correct closeout move is therefore:
\begin{enumerate}
\item freeze the strongest exact reference window already proven in canonical docs,
\item hash the committed code and documentation inputs,
\item name every downstream runtime artifact location precisely,
\item and record which stages still lack materialized outputs.
\end{enumerate}

That converts a vague future workflow into a precise reproducibility contract. When Task 2 is revisited with real runtime outputs present, the same manifest path can be regenerated and upgraded from \texttt{frozen\_definition\_with\_missing\_runtime\_outputs} to a materially complete end-to-end evidence packet.

\end{document}
